% Source: https://artofproblemsolving.com/community/c3h1568598_a_little_argument_on_probability_of_getting_hit_by_a_car
% Stable identity: aops-topic-1568598-q01
Let $T$ be the event ``hit today'' and $M$ the event ``hit tomorrow''.

\textbf{(a)} By independence, $P(T\cap M)=P(T)P(M)=\boxed{p^2}$.

\textbf{(b)} By the definition of conditional probability,
\[
P(M\mid T)=\dfrac{P(T\cap M)}{P(T)}=\dfrac{p^2}{p}=\boxed{p}.
\]
This is exactly what independence means: knowing that $T$ happened does not change the probability of $M$.

\textbf{The friend's error.} The number $p^2$ is the probability, judged \emph{before today}, that both events happen. Once we know that the person was hit today, the event $T$ is certain, and only tomorrow is uncertain, with probability $p$. The friend confuses $P(T\cap M)$ with $P(M\mid T)$. (Whether real accidents on different days are independent is a question about the model, not about probability theory: if being hit made a person more careful, the events would not be independent and the model would have to change.)
